\documentclass[12pt]{article}
\usepackage[utf8]{inputenc}
\usepackage[T1]{fontenc}
\usepackage{hyperref}
\title{Wathawurrung Language Resources for K-12 Curriculum Development}
\author{Compiled by Claude Sonnet 5.5 (claude-sonnet-5-5) for Andr\'as Kornai}
\date{}
\begin{document}
\maketitle

\section{Introduction}
Wathawurrung (also Wadawurrung, Wathaurong; ISO 639-3 wth) is a language of the Geelong and Ballarat region of Victoria, Australia. Glottolog classifies it as extinct, but revitalization work is under way, notably through the Victorian Aboriginal Corporation for Languages (VACL).

\section{Primary Educational Resources}
VACL runs a Schools Digital Resources Project aimed at audio-visual resources for language reclamation in Victorian schools. I did not find a Wathawurrung-specific school programme in the results.
\begin{itemize}
\item VACL Schools Digital Resources Project: \url{https://vacl.org.au/project/schools-digital-resources-project/}
\end{itemize}

\section{Linguistic and Reference Materials}
\begin{itemize}
\item VACL Wathaurong / Wadda Wurrung page. It mentions the ``Learning Wathaurong'' CD-ROM, the Wathawoorrong Dictionary, a list of about 7,000 words, and an iPad app: \url{https://vacl.org.au/language/wathaurong-wadda-wurrung/}
\item Blake (ed.), \emph{Wathawurrung and the Colac language of southern Victoria} (Pacific Linguistics, 1998), open access via the ANU repository: \url{https://openresearch-repository.anu.edu.au/bitstreams/098c3e70-6d29-4f74-b631-7cec113d2e9a/download}
\item Wikipedia: \url{https://en.wikipedia.org/wiki/Wadawurrung_language}
\item Glottolog: \url{https://glottolog.org/resource/languoid/id/wath1238}
\end{itemize}

\section{Religious and Community Resources}
\begin{itemize}
\item Inspiring projects around Australia (University of Melbourne Research Unit for Indigenous Language): \url{https://arts.unimelb.edu.au/research-unit-for-indigenous-language/resources/inspiring-projects}
\end{itemize}

\section{Digital and Audio Resources}
VACL's page mentions an iPad app and a CD-ROM, but I did not see direct links to them.

\section{Curriculum Development Context}
The language is documented enough for revitalization work: a dictionary, a word list of about 7,000 words, and an academic volume. I found no current contact or availability details for the dictionary. VACL is the natural partner for any school curriculum, working with the Wadawurrung community.
\end{document}